% Tu izpolnite podatke vašega doktorskega dela
\newcommand{\naslovSlo}{Naslov dela}
\newcommand{\naslovAng}{Thesis title}
\newcommand{\avtor}{Ime Priimek avtorja}

\newcommand{\tipDela}{Doktorsko delo}
% Izberite obliko glede na spol kandidata oz. kandidatke.
\newcommand{\predlozil}{Predložil}
% \newcommand{\predlozil}{Predložila}

\newcommand{\mentor}{akad.~naziv Ime Priimek, izobrazba}
\newcommand{\somentor}{akad.~naziv Ime Priimek, izobrazba} % Prazno, če somentorja ni.

\newcommand{\zagovorMesec}{mesec}
\newcommand{\zagovorLeto}{leto zagovora}

% Podatki za načrt ravnanja z raziskovalnimi podatki (NRRP)
\newcommand{\doktorskiProgram}{Doktorski študijski program Strojništvo, področje \ldots}

% Podatki za generacijo strani Izvleček/Abstract
\newcommand{\UDK}{XXX.XX:XXX.(XXX.X)}

\newcommand{\tekocaStevilka}{DR III/XX} % Prepišite z odločbe o potrjeni temi.

% Navedite 5--8 ključnih besed, vsako v svoji vrstici. Slovenske ključne
% besede praviloma pišemo z malo začetnico.
\newcommand{\kljucneBesede}{%
ključna beseda 1 (vse ključne besede pišemo z malo začetnico)\\
ključna beseda 2 (zahtevanih je 5--8 ključnih besed)\\
ključna beseda 3\\
ključna beseda 4\\
ključna beseda 5\\
ključna beseda 6}

\newcommand{\izvlecek}{% Nadomestite z dejanskim izvlečkom (praviloma 150--300 besed, do pol strani).
V začetku izvlečka je jedrnat opis obravnavanega problema. Sledi opis izbrane metodologije in nato opis najpomembnejših rezultatov oz. ugotovitev brez dodatnih pojasnjevanj. Obseg izvlečka je praviloma omejen na 150 do 300 besed oziroma naj ne presega polovice strani. Študent pridobi UDK številko v knjižnici FS, za katero mora zaprositi vsaj 3 delovne dni pred oddajo. To lahko naredi na podlagi navedenih ključnih besed in povzetka, ki jih mora študent poslati v knjižnico na naslov knjiznica@fs.uni-lj.si. Kandidat mora vnesti tudi tekočo številko doktorskega dela, naslov dela, svoje ime in priimek ter ključne besede.}

\newcommand{\keywords}{%
keyword 1 (write all keywords in lower case)\\
keyword 2 (five to eight keywords are required)\\
keyword 3\\
keyword 4\\
keyword 5\\
keyword 6}

\newcommand{\abstract}{%
At the beginning of the abstract, a concise description of the addressed problem must be provided. This is followed by a description of the chosen methodology and then by a presentation of the most important results or findings, without additional explanations. The abstract should generally be limited to 150 to 300 words and should not exceed half a page. The candidate must enter the UDC number, which is obtained from the library based on the submitted information about the thesis: the title, the supervisor and any co-supervisors, at least five keywords, and the abstract. The candidate must also enter the serial number of the thesis, the title of the work, their name and surname, and the keywords. Keywords are generally written in the plural.}

\newcommand{\datotekaTeme}{tema.pdf} % Odločba o potrjeni temi; vključijo se vse strani PDF.

\newif\ifzahvala

\zahvalatrue % Za izpust zahvale uporabite \zahvalafalse.
